\section{Reusable Temporal Assignment-Map Architecture}
\label{sec:planned-temporal-assignment-map}

This section presents the architecture implemented by the current package and
identifies the deliberately incomplete extensions.  It extracts a stable
computational interface from the forward models developed earlier in this
part, especially the scheduled time-expanded formulation, the reduced
journey-response formulation, the additive fixed-routing decomposition, and
the measurement model.  It deliberately distinguishes verified
public-library functionality from proposed work.  In particular, it does not
specify an alternative supply model: the authoritative supply backend remains
scheduled, time-expanded assignment, while estimation uses its persisted
fixed-routing response.

The following labels are used throughout this section.
\begin{description}
\item[\textbf{IMPLEMENTED}] The public library contains a concrete, tested
  implementation of the stated behavior for the current canonical
  direct-scheduled workflow.  Broader generalization is stated explicitly
  where it remains incomplete.
\item[\textbf{PARTIALLY IMPLEMENTED}] Relevant tested foundations exist, but
  the complete contract or end-to-end workflow specified here does not.
\item[\textbf{EXPERIMENTAL}] A callable prototype exists, but its interface or
  operational guarantees are not yet a stable general contract.
\item[\textbf{UNIMPLEMENTED}] The component is a proposed design and is not
  currently supplied by the public library.
\end{description}
When a component spans several existing modules, the label refers to the
complete component described here, not to every useful foundation beneath it.

\subsection{Layered forward decomposition}

Let $\theta\in\mathbb{R}^{p}$ be a low-dimensional vector of demand-model
parameters.  A demand generator $g$ produces the complete physical OD-time
demand vector $x\in\mathbb{R}^{n}$ in a prescribed canonical order.  A
preprocessed assignment map $A\in\mathbb{R}^{m\times n}$ maps that demand to
predicted boarding and alighting measurements
$\widehat y\in\mathbb{R}^{m}$:
\begin{equation}
  \theta
  \overset{g}{\longmapsto}
  x
  \overset{A}{\longmapsto}
  \widehat y,
  \qquad
  \widehat y=A\,g(\theta).
  \label{eq:planned-layered-forward-map}
\end{equation}
Assignment, demand, observation, and estimation are independent layers.  The
assignment layer defines how passenger demand creates expected measurement
events.  The demand layer defines $g$.  The observation layer compares
$\widehat y$ with recorded counts.  The estimation layer chooses an algorithm
for varying $\theta$.  This separation makes the expensive supply-dependent
calculation reusable across model specifications.

\paragraph{Implementation status.}
\textbf{IMPLEMENTED for the current direct-scheduled full-network workflow.}
The public library composes gravity demand generation, measurement operators,
likelihoods, and optimizers in this form.  The content-addressed full-network
temporal assignment artifact has been constructed, persisted, loaded in a
fresh process, and benchmarked through the same forward/adjoint contract.
Generalization to every possible demand, observation, and assignment backend
remains a separate roadmap.

\subsection{Strict assignment-map contract}
\label{sec:strict-assignment-map-contract}

The external numerical responsibility of an assignment component is limited
to the two products
\begin{equation}
  x\longmapsto Ax,
  \qquad
  r\longmapsto A^{\mathsf T}r,
  \label{eq:assignment-forward-adjoint-contract}
\end{equation}
where $r\in\mathbb{R}^{m}$ is a measurement-space vector.  Neither product
requires a dense realization of $A$.

An assignment artifact may depend on network topology, the scheduled
timetable, temporal discretization, fixed route-choice parameters, fixed
departure-time-choice parameters, transfer, waiting, and journey-duration
constraints, measurement definitions, and coefficient-pruning tolerances.
These inputs determine the physical response represented by $A$.

It must not depend on a gravity-model specification, gravity parameters,
priors, the likelihood, the optimizer, or estimated OD values.  Consequently,
changing a demand model, prior, likelihood, or optimizer must not trigger
assignment preprocessing.  Changing the timetable, temporal intervals,
route-choice parameters, departure-choice parameters, or feasibility
constraints must invalidate the artifact.  Measurement-definition changes
must likewise invalidate the corresponding measurement map.

\paragraph{Implementation status.}
\textbf{IMPLEMENTED for the direct-scheduled artifact path.}  Tested
linear-operator and gravity-measurement protocols expose forward and adjoint
products, shapes, fingerprints, and operational capabilities.  The persisted
artifact identity separates physical indexing, supply, timetable, temporal,
choice, feasibility, measurement, coefficient-policy, numerical, and schema
dependencies, with explicit mismatch diagnostics.  Existing fixed-routing
operators validate against this identity through an adapter.  The broader
dependency boundary is not yet shared by every assignment backend and demand
model.

\subsection{Assignment-backend abstraction}

Downstream code should depend on a stable abstract interface rather than on the
internal time-expanded representation.  A minimal conceptual contract is:
\begin{verbatim}
class AssignmentOperator:
    number_of_demand_cells: int
    number_of_measurements: int

    def matvec(self, demand):
        """Calculate predicted measurements A @ demand."""

    def rmatvec(self, residual):
        """Calculate the adjoint product A.T @ residual."""
\end{verbatim}
Production implementations additionally need a numerical type, immutable
identity metadata, and validated dimensions.  Those additions must not weaken
the two-operation numerical contract in
Equation~\eqref{eq:assignment-forward-adjoint-contract}.

Other assignment backends could later implement the same canonical indexing
and forward/adjoint contract.  Such a substitution must not require changes to
demand models, observation models, likelihoods, or optimizers.

\paragraph{Implementation status.}
\textbf{IMPLEMENTED.}  The public library provides the runtime-checkable
\texttt{AssignmentOperator} protocol with canonical dimensions, index and
artifact fingerprints, and the required \texttt{matvec} and \texttt{rmatvec}
products.  A validated adapter exposes existing explicit, matrix-free, and
sharded fixed-routing operators through this contract without altering their
numerical products.  The authoritative scheduled time-expanded calculation is
also available through the same contract as a deliberately unoptimized
reference backend.

\subsection{Canonical indexing and compatibility}

Three stable canonical tables define the spaces on which $A$ operates:
\begin{enumerate}
\item a physical OD--departure-interval cell table, with one immutable row for
  every demand coordinate;
\item a boarding and alighting measurement table, with one immutable row for
  every predicted count; and
\item a temporal-interval table giving ordered interval identifiers and exact
  boundaries.
\end{enumerate}
Every demand model must return values in the assignment artifact's canonical
OD-time order.  Display labels are not identities: canonical rows must retain
stable physical identifiers and interval identifiers.  Compatibility checks
must cover dimensions, schemas, and fingerprints and must fail explicitly,
rather than silently reordering, truncating, or accepting incompatible inputs.

\paragraph{Implementation status.}
\textbf{IMPLEMENTED for the direct-scheduled artifact path.}  Existing
reduced-OD and fixed-routing paths use immutable layouts, explicit dimensions,
fingerprints, schemas, and fail-closed cache validation.  Immutable canonical
interval, physical OD-time, and boarding/alighting measurement tables are now
populated by the direct-scheduled builder and persisted in the artifact
manifest, including separate physical-artifact and free/fixed-binding
fingerprints.  They are not yet populated and consumed by every assignment
backend and every demand-model workflow.

\subsection{Non-uniform temporal discretization}

The proposed temporal partition consists of ordered, non-overlapping intervals
\begin{equation}
  I_1,I_2,\ldots,I_K.
\end{equation}
Intervals should be finer when supply, demand, route availability, or transfer
conditions change rapidly, and broader during stable conditions.  Selecting
the partition is therefore a model-design problem, not merely a global choice
of interval length.  A candidate partition should balance
\begin{equation}
  \text{temporal approximation loss}
  +\lambda\,\text{assignment complexity},
  \label{eq:temporal-partition-objective}
\end{equation}
where $\lambda\geq 0$ expresses the desired precision--resource compromise.

Candidate diagnostics include service frequencies and headways, active lines
and route patterns, running times, transfer opportunities, boarding and
alighting count profiles, approximate assignment support, and projected
nonzeros, memory, and execution time.  Rather than impose one arbitrary
partition, preprocessing should eventually report a precision--complexity
Pareto frontier from which a documented choice can be made.

\paragraph{Implementation status.}
\textbf{PARTIALLY IMPLEMENTED.}  The public preprocessing module provides
\texttt{recommend\_time\_discretization} and its CSV convenience wrapper.  It
scores explicit candidate partitions from timestamped count data and returns a
documented recommendation together with candidate diagnostics.  The selected
intervals must still be reviewed and materialized by the case owner before
they become the case input.  There is no universal automatic partition
selector or Pareto-frontier search integrated into every case-study workflow.

\subsection{Scheduled time-expanded preprocessing}

The scheduled time-expanded graph described earlier in this part remains the
authoritative routing representation during preprocessing.  For each physical
OD pair and departure interval $s$, preprocessing should:
\begin{enumerate}
\item construct feasible joint departure-time and route alternatives;
\item apply fixed exogenous departure-time-choice and route-choice parameters;
\item calculate the alternative probabilities;
\item project boarding and alighting contributions onto measurement interval
  $t$; and
\item consolidate those contributions into a sparse temporal block $A_{t,s}$.
\end{enumerate}
Departure-time and route choice are thus combined inside scheduled
time-expanded preprocessing.  If $x_s$ denotes the demand coordinates in
departure interval $s$ and $y_t$ the predicted measurements in interval $t$,
the resulting operator satisfies
\begin{equation}
  y_t=\sum_s A_{t,s}x_s.
  \label{eq:temporal-block-forward-map}
\end{equation}
The time-expanded graph is not traversed during estimation; repeated objective
and gradient evaluations use only the persisted assignment response.

\paragraph{Implementation status: scheduled routing foundation.}
\textbf{IMPLEMENTED.}  The public library contains scheduled time-expanded
network construction and assignment foundations, including journey and
measurement-response preprocessing.  A canonical scheduled reference operator
now evaluates the complete time-expanded forward calculation and obtains its
exact adjoint by reverse-mode differentiation.  It validates canonical and
supply-dependent identities and is intended for small equivalence tests rather
than repeated production estimation.

\paragraph{Implementation status: fixed-routing linear operations.}
\textbf{IMPLEMENTED.}  Dense, sparse, matrix-free, sharded, and selected-block
fixed-routing paths provide tested forward products; the relevant estimation
operators also provide adjoint products and bounded operational controls.

\paragraph{Implementation status: complete temporal-block conversion.}
\textbf{IMPLEMENTED for the current canonical direct-scheduled path.}  The
library can derive an exact, independently
addressable sparse family $\{A_{t,s}\}$ column by column from the authoritative
scheduled reference operator, with progress, ETA, deadline handling, duplicate
consolidation, and coefficient-mass diagnostics.  A faster constructor batches
unit responses in one JIT-compiled fixed-shape kernel, pads the final chunk, and
aggregates measurements on device before transferring bounded response chunks.
It compiles once and never materializes the complete dense map.  The production
constructor now goes further: it prepares fixed scheduled routing probabilities
once, propagates only measurement-supported edges, emits sparse measurement
contributions into deterministic shards, checkpoints them atomically, resumes
after interruption, selectively rebuilds corrupt shards, and then finalizes
canonical temporal blocks.  A general explicit joint departure-time-choice
layer beyond the current canonical departure intervals remains unfinished.
The resulting block operator is exposed through the established gravity
measurement-operator protocol, including forward, adjoint, and batched products,
so gravity objectives require no representation-specific changes.

\subsection{Sparse temporal block structure}

The intended block map is sparse for three distinct reasons.  First,
\emph{causality} prevents demand from affecting earlier measurements.  Second,
\emph{temporal banding} follows from a maximum feasible journey duration, which
limits the active range of $t-s$.  Third, \emph{spatial sparsity} arises because
one OD journey visits only a small subset of all measured locations and event
types.

Construction of each block should merge duplicate row--column contributions,
remove exact zeros, and optionally prune negligible probabilities using an
explicit, auditable tolerance.  Pruning must retain diagnostics for removed
probability mass and its distribution.  The representation should use compact
integer indices, persist temporal blocks independently, and support efficient
forward and adjoint access.  Compressed sparse row (CSR) storage is natural for
$Ax$; compressed sparse column (CSC), or equivalently a persisted transposed
CSR representation, avoids repeated conversion for $A^{\mathsf T}r$.

\paragraph{Implementation status.}
\textbf{IMPLEMENTED for the persisted temporal-block path.}  Exact in-memory temporal blocks now merge duplicates,
remove zeros, retain coefficient-mass diagnostics, and provide JAX-compatible
forward and adjoint products.  A versioned manifest persists blocks
independently and validates their content hashes before reuse.  A CPU backend
constructs and retains CSR for forward products and CSC for adjoint products;
its custom JAX reverse rule returns the CSC adjoint without materializing a
dense matrix.  Direct alternative-level construction and the complete
probability-pruning workflow remain unfinished.

\subsection{Block-support complexity profiling}

Before expensive construction, a bounded profiler should estimate active block
support, projected nonzeros, persistent storage, peak construction memory, and
forward/adjoint execution cost.  Its estimates should be reported by temporal
block and in aggregate so that an infeasible partition can be rejected before
full preprocessing.  The profiler is also the computational-complexity input
to the Pareto analysis in
Equation~\eqref{eq:temporal-partition-objective}.

\paragraph{Implementation status.}
\textbf{PARTIALLY IMPLEMENTED.}  Fixed-routing block-coordinate workflows
provide bounded, resumable support preflight, resource limits, retained-state
limits, progress, and checkpoint validation.  A canonical temporal profiler now
uses causality and a maximum journey duration to report feasible blocks,
candidate entries, excluded entries, and conservative storage.  It does not
yet estimate spatial support, peak construction memory, or execution time from
the scheduled graph, nor compare candidate temporal partitions.

\subsection{Independent demand-model interface}

The demand layer should depend only on its parameters, explanatory data, and
the canonical demand-cell table.  A conceptual differentiable interface is:
\begin{verbatim}
class DemandModel:
    parameter_names: tuple[str, ...]

    def demand(self, parameters):
        """Generate demand in canonical OD-time order."""

    def transpose_jacobian_product(
        self, parameters, demand_gradient
    ):
        """Calculate J_g(parameters).T @ demand_gradient."""
\end{verbatim}
Possible implementations include minimal global gravity, time-specific
gravity, origin-group production, destination-group attraction, low-rank OD
interactions, gravity plus residual corrections, directly parameterized OD
demand, and externally supplied demand.  Switching between these models must
reuse the same compatible assignment artifact.

\paragraph{Implementation status.}
\textbf{PARTIALLY IMPLEMENTED.}  The generic reduced-demand code supplies
composable production, attraction, impedance, and low-rank interaction
specifications; canonical demand generation is differentiable and several
progressive gravity specifications are reusable with the same response
operator.  Existing fixed, directly supplied, and reduced parameterizations
provide additional foundations.  There is not yet one stable public
\texttt{DemandModel} protocol covering every listed family, explicit canonical
compatibility, and a named transpose-Jacobian product.

\subsection{Independent observation model and differentiation}

The observation layer receives observed counts and $Ax$, not routing
structures.  It should select independently among Poisson, negative-binomial,
possible zero-inflated formulations, and future count likelihoods.  With
\begin{equation}
  x=g(\theta),
  \qquad
  \widehat y=Ax,
\end{equation}
the chain rule gives
\begin{equation}
  \nabla_\theta L
  =J_g(\theta)^{\mathsf T}
   A^{\mathsf T}
   \nabla_{\widehat y}L.
  \label{eq:layered-adjoint-gradient}
\end{equation}
Thus neither the dense assignment matrix nor the full gravity Jacobian needs
to be materialized.  The observation component controls the distributional
assumption, dispersion, detection terms when modeled, zero-inflation design,
and calibration mask, while remaining independent of assignment internals.

\paragraph{Implementation status.}
\textbf{PARTIALLY IMPLEMENTED.}  The current gravity objective implements
Poisson and negative-binomial count likelihoods, and automatic differentiation
composes either one with the fixed-routing operator products.  The aggregate
GPS channel provides separate multinomial, Dirichlet--multinomial, and
tempered-multinomial contributions when enabled.  Zero-inflated count forms
remain future work, and a stable observation-model protocol shared by every
estimation path is not yet defined.

\subsection{Reusable assignment artifact}

The conceptual persisted layout is:
\begin{verbatim}
assignment_map/
|-- manifest.json
|-- fixed_measurement_offset.npy
`-- blocks/
    `-- block-*.npz
\end{verbatim}
The current artifact embeds the canonical OD-time, measurement, and temporal
interval tables in the manifest under \texttt{canonical\_index}.  The manifest also
records a schema version, canonical-index fingerprints, network and timetable
fingerprints, temporal discretization, choice-parameter fingerprints,
feasibility constraints, coefficient precision and threshold, dimensions,
total nonzeros, block dimensions, content hashes, and construction
diagnostics.  It does not contain a gravity specification, prior, likelihood,
or optimizer configuration.  Run-level progress, checkpoint, and benchmark
records are maintained alongside the artifact rather than as required payload
subdirectories.  Payloads are published atomically and validated before reuse;
incompatible or incomplete artifacts fail closed.

\paragraph{Implementation status.}
\textbf{IMPLEMENTED for the current direct-scheduled artifact.}  The library
has a content-addressed,
versioned temporal-block artifact with an atomic publication step, canonical
tables, complete identity metadata, per-block content hashes, fixed-offset
validation, and fail-closed fresh-process loading.  Its direct construction
workspace contains deterministic validated shards and a checkpointed manifest;
completed shards survive interruption and invalid shards are rebuilt
selectively before atomic final publication.  The full-network artifact was
loaded in a fresh process and used for a measured forward product; diagnostic,
progress, and benchmark records remain run-level products rather than required
contents of the immutable assignment artifact.

Production activation is policy controlled.  Explicit \texttt{off} leaves the
existing loader unchanged, while explicit \texttt{direct} constructs or resumes
the block artifact.  In \texttt{auto} mode, a valid persistent artifact is
reused before preparing routing; otherwise construction occurs only when the
expected evaluation savings strictly exceed its estimated construction cost.
Invalid final artifacts are rejected, moved to a uniquely named quarantine
directory for diagnosis, and rebuilt through the atomic publication path.

\subsection{Validation requirements}

The scheduled temporal assignment map must pass the following validation
sequence before it can replace the authoritative preprocessing calculation in
estimation:
\begin{enumerate}
\item compare scheduled sparse blocks with the authoritative time-expanded
  assignment;
\item test forward equivalence on several deterministic random vectors;
\item verify the adjoint identity
  \begin{equation}
    \langle Ax,r\rangle
    =\langle x,A^{\mathsf T}r\rangle;
    \label{eq:planned-adjoint-validation}
  \end{equation}
\item measure nonzeros, storage, peak memory, and forward/adjoint execution
  time;
\item report approximation errors by time interval, transfer count, and
  journey class;
\item track affected observed count mass, not only matrix norms;
\item verify artifact invalidation and canonical-index compatibility; and
\item verify that changing only the demand model reuses the same assignment
  artifact.
\end{enumerate}
Tolerance choices, random seeds, coefficient thresholds, and reference
fingerprints belong in the validation record.  Exact equivalence tests should
be separated from explicitly approximate pruning tests.

\paragraph{Implementation status.}
\textbf{PARTIALLY IMPLEMENTED.}  Existing operator tests cover forward
equivalence, adjoint products, cache incompatibility, fingerprints, numerical
agreement, frozen-demand offsets, construction progress, interruption and
resume, finalized-artifact reuse, selective shard repair, and selected
performance diagnostics.  Production-policy tests additionally cover explicit
disablement, economically unjustified automatic decline without routing work,
fresh-process reuse, gravity-protocol matrix products, and safe quarantine and
rebuilding of an incompatible final artifact.  Stratified approximation and
affected-count-mass reporting remain to be implemented for this representation.

The executable public validation
\texttt{docs/source/examples/direct\_scheduled\_gravity\_validation.py} applies
this interface end to end to \texttt{simple\_example\_02}.  It constructs the
canonical demand and observation intervals, reports shard progress, evaluates
the same negative-binomial gravity objective and adjoint gradient through both
the established BCOO operator and the temporal-block operator, and supports a
mandatory fresh-process cache-reuse check.  On the development laptop, the
first validation constructed 1,201 nonzeros in 0.31 seconds (0.48 seconds for
complete activation) and stored 15,492 bytes.  The objective absolute difference
was $3.05\,10^{-5}$ and the largest gradient absolute difference was
$1.91\,10^{-6}$.  A second Python process loaded the persistent artifact in
0.011 seconds, performed no construction, and reproduced the same numerical
differences.  The test suite executes this two-process protocol automatically.

Construction now uses a single monotonic deadline contract shared by activation,
routing guards, support and planning guards, shard validation and construction,
temporal assembly, and atomic persistence.  A normal bounded stop returns a
structured terminal record rather than a partial operator or an undifferentiated
error.  The record identifies the phase, elapsed and remaining time, completed
work, next resumable unit, checkpoint and artifact locations, predicted next
cost, checkpoint reusability, and any overshoot caused by an indivisible
operation.  Versioned progress events use the same phase vocabulary and remain
independent of presentation libraries.

Shard construction stops only after atomically committing and validating the
current shard, and a later invocation reuses it.  Fixed-shape destination-group
routing batches,
destination-group origin-support fragments, aggregate materialized support, the
deterministic plan, and one temporal fragment per numerical source shard are
also content validated and persisted.
Thus resume bypasses the routing factory, support discovery, and planning when
their checkpoints are compatible, and temporal assembly restarts at its first
missing fragment.  Final persistence checks the deadline between blocks and
removes an unpublished staging directory on termination.  Valid
completed artifacts retain cache-first priority even under an almost-expired
new budget. Routing is compiled once for a bounded padded batch shape; each
batch is synchronized, transferred, content validated, and atomically persisted
before the next dispatch. Downstream support and numerical construction load at
most one routing batch and never concatenate the global destination-by-link
arrays. JAX compilation and one batch execution remain indivisible internally:
they are guarded before launch and record overshoot. Origin-support discovery checks
the deadline between bounded origin chunks and atomically commits every
completed destination group, so resume continues at the first missing group.
Numerical measurement-shard construction may now use an explicitly bounded
thread pool.  All workers share the one compiled executable, but each owns its
own routing-batch reader.  Per-worker workspace preflight and a resident-shard
ceiling bound admission.  Workers write validated payloads only to isolated
staging locations; the parent publishes them and advances the manifest in
canonical plan order, even when workers finish out of order.  Worker count is
excluded from scientific provenance, interrupted parallel and serial runs
therefore reuse the same cache, and failures remove unpublished staging files
without invalidating previously committed shards.
Origin-specific support materialization has a separate positive configurable
entry ceiling, whose default is 125 million.  This ceiling is an operational
memory guard rather than scientific provenance, so increasing it preserves
compatible destination-group support checkpoints.  Summary-only analysis can
avoid materialization, but production numerical construction cannot.
Numerical-shard preflight diagnostics likewise retain actual and permitted
values for shard count, manifest size, filesystem work, sparse products,
construction dispatches, and worker memory.  Unsafe plans raise a structured
subclass of \texttt{MemoryError}.  Shard-sizing controls remain provenance
because they alter persisted identities; purely operational ceilings are stored
in the plan and re-evaluated at reuse time, so an evidence-based ceiling change
does not discard compatible support checkpoints.
Only the current CPU chunk or destination group may be repeated. Therefore the
deadline hardening is sufficient for bounded public full-network probes.

\subsection{Implementation roadmap}

Table~\ref{tab:temporal-assignment-roadmap} records the audited status of the
complete components, using the conservative convention at the start of this
section.  Priorities order the remaining architectural work; a dash indicates
an existing foundation rather than a new implementation priority.

\begin{table}[htbp]
\centering
\footnotesize
\caption{Roadmap for the reusable temporal assignment-map architecture.}
\label{tab:temporal-assignment-roadmap}
\begin{tabular}{|p{0.46\textwidth}|p{0.29\textwidth}|>{\centering\arraybackslash}p{0.09\textwidth}|}
\hline
\textbf{Component} & \textbf{Status} & \textbf{Priority} \\
\hline
Scheduled time-expanded assignment foundation & IMPLEMENTED & -- \\
\hline
Scheduled time-expanded reference operator & IMPLEMENTED & -- \\
\hline
Fixed-routing linear operations & IMPLEMENTED & -- \\
\hline
Stable assignment-operator abstraction & IMPLEMENTED & -- \\
\hline
Canonical physical OD-time index & IMPLEMENTED & -- \\
\hline
Data-driven non-uniform partition recommendation & PARTIALLY IMPLEMENTED & 1 \\
\hline
Block-support complexity profiler & PARTIALLY IMPLEMENTED & 1 \\
\hline
Scheduled $A_{t,s}$ construction & IMPLEMENTED & -- \\
\hline
Sparse temporal-block persistence & IMPLEMENTED & -- \\
\hline
Optimized temporal CSR/CSC forward and adjoint & IMPLEMENTED & -- \\
\hline
Independent demand-model contract & PARTIALLY IMPLEMENTED & 2 \\
\hline
Interchangeable gravity specifications & IMPLEMENTED & -- \\
\hline
Independent observation models & PARTIALLY IMPLEMENTED & 3 \\
\hline
Assignment-backend substitution test & UNIMPLEMENTED & 4 \\
\hline
\end{tabular}
\end{table}

\paragraph{Implementation status.}
\textbf{IMPLEMENTED for the current canonical direct-scheduled workflow.}  The
full-network path now includes the canonical index, fixed scheduled routing,
temporal sparse-block construction, content-addressed persistence, atomic
resumption, and the tested CSR/CSC forward and adjoint backend.  The remaining
priorities are generalizations: data-driven partition selection, broader
demand and observation protocols, backend substitution, direct alternative-level
construction, and a more general joint departure-time-choice layer.

\subsection{Architectural invariants}

The following requirements are non-negotiable for implementation of this
roadmap:
\begin{enumerate}
\item \emph{Assignment preprocessing is independent of the gravity
  specification.}
\item \emph{Demand models produce vectors in canonical OD-time order.}
\item \emph{Assignment operators provide forward and adjoint operations.}
\item \emph{Estimation never traverses the time-expanded graph.}
\item \emph{Changing only the demand model does not rebuild the assignment
  artifact.}
\item \emph{Approximation and pruning are accompanied by quantitative
  diagnostics.}
\item \emph{Long preprocessing operations are resumable and report progress.}
\item \emph{Downstream estimation code depends on the assignment contract, not
  on the scheduled time-expanded implementation.}
\end{enumerate}
The eighth invariant is the sole future-proofing requirement for possible
alternative assignment representations; no additional supply methodology is
specified here.

For public preprocessing, progress is a machine-readable contract rather than
a presentation-layer convention. Timetable indexing, measurement-response
aggregation, origin-support discovery, structural-zero classification, and
artifact publication report completed and total units, elapsed time, observed
throughput, estimated remaining time, confidence, and the current unit. Events
occur at phase boundaries, approximately every one percent, or at the first
unit boundary after ten seconds. The ETA is explicitly unavailable before the
first measured unit. A library or backend call that cannot expose internal
work remains an identified indivisible unit and is not assigned a fictitious
completion estimate.

\paragraph{Implementation status.}
\textbf{IMPLEMENTED for the current direct-scheduled full-network workflow.}
The current artifact path collectively enforces the canonical contract, index
tables, forward/adjoint products, demand-model separation, progress reporting,
resumability, and fail-closed artifact reuse.  Approximation-specific
affected-count-mass reporting and substitution across independent assignment
backends, demand protocols, and observation protocols remain future work.
